\chapter{Aplicație: EEG}
\label{ch:eeg}

Activitatea electrică a creierului, culeasă cu electrozi pe scalp, se
citește pe benzi de frecvență cu semnificație clinică proprie: delta
(0.5-4~Hz, somn profund), theta (4-8~Hz, relaxare adâncă), alfa (8-13~Hz,
ochii închiși, relaxare trează), beta (13-30~Hz, concentrare activă) și
gamma (30-80~Hz, procesare cognitivă complexă). Polisomnografia (studiul
somnului), diagnosticul de epilepsie și interfețele creier-calculator se
bazează, în forme diferite, pe separarea acestor benzi, un ceas care
„estimează” somnul din mișcare și puls nu are acces la niciuna dintre ele.
EEG-ul adevărat rămâne un instrument de laborator sau de clinică.

Demo-ul din tur (Figura~\ref{fig:eeg-demo}) extrage o singură bandă
dintr-un semnal EEG sintetic cu un filtru trece-bandă obișnuit
(Secțiunea~\ref{sec:filtre}), nu cu transformata wavelet discretă folosită
în restul documentului. Motivul e scris direct pe ecran: DWT dublează
frecvența la fiecare nivel, o octavă, dar granițele benzilor clinice nu sunt
spațiate octavă-cu-octavă, delta-theta e aproape un factor de doi,
delta-alfa e deja un factor de $3.25$, așa că niciun set fix de niveluri
DWT nu s-ar alinia cu toate granițele deodată. Un filtru trece-bandă clasic,
cu tăieturi alese liber, nu are această constrângere.

\begin{figure}[htbp]
\centering
\includegraphics[width=\textwidth]{37-applications-eeg-demo.png}
\caption{Banda alfa (8-13~Hz) extrasă cu un filtru trece-bandă dintr-un EEG
  sintetic. Nota din josul figurii explică de ce nu DWT.}
\label{fig:eeg-demo}
\end{figure}
